% TOP_PROBLEM: {"id": 30007374, "title": "The thermodynamic pressure of every bounded Gaussian Ising perceptron activation", "category_id": 19, "category": {"id": 19, "name": "probability", "display_name": "Probability", "description": "Problems involving probability theory, stochastic processes, and random structures.", "slug": "probability", "order_index": 19, "created_at": "2026-07-31T15:26:25.671Z"}, "status": "open"}
% FIRST_POSTED: 2026-10-06
% ATTEMPT_STATUS: solved
% OPENAI_CLAIM: {"schema_version":1,"source_repo":"https://github.com/openai/math","source_commit":"adc7f1241b42e322a6451854ab7e4b4c146bf78a","release_date":"2026-10-06","match":"full","resolution":"proved","independently_reviewed":false,"second_pass":"agreed","families":[{"family":"222","title":"Perceptron free energies and microscopic jamming exponents","subject":"Probability and statistical mechanics","manuscripts":[{"title":"The free energy of the Ising random perceptron","pdf_path":"preprints/The-free-energy-of-the-Ising-random-perceptron-September-24-2026/The-free-energy-of-the-Ising-random-perceptron-September-24-2026.pdf","readme_path":"preprints/The-free-energy-of-the-Ising-random-perceptron-September-24-2026/README.md","date":"2026-09-24","theorem_ref":"preprints/The-free-energy-of-the-Ising-random-perceptron-September-24-2026/build/sections/model.tex:88 (thm:main)"}],"lean":{"doc_path":"lean/docs/222.md","comparators":["lean/ComparatorChallenges/IsingFiniteness.json","lean/ComparatorChallenges/IsingFiniteness.lean","lean/ComparatorChallenges/PerceptronFreeEnergy.json","lean/ComparatorChallenges/PerceptronFreeEnergy.lean","lean/ComparatorChallenges/SphericalField.json","lean/ComparatorChallenges/SphericalField.lean"],"declarations":[],"covers_main_theorem":false},"reasoning_trace":null}]}
\documentclass[11pt]{article}
\usepackage[margin=25mm]{geometry}
\usepackage{fontspec}
\setmainfont{Latin Modern Roman}
\usepackage{amsmath,amssymb,mathtools}
\usepackage{unicode-math}
\setmathfont{Latin Modern Math}
\usepackage{enumitem}
\usepackage{xurl,hyperref,bookmark}
\hypersetup{colorlinks=true,urlcolor=blue}
\setcounter{secnumdepth}{0}
\setlength{\parindent}{0pt}
\setlength{\parskip}{5pt}
\setlength{\emergencystretch}{3em}
\begin{document}
\section*{The thermodynamic pressure of every bounded Gaussian Ising perceptron activation}

\subsection{OpenAI's claim}
OpenAI presents a proof in the selected manuscript. The Ising Gaussian perceptron is claimed to have an exact pressure formula for every bounded Borel log-potential at positive temperature and density.

\subsection{Outline of the argument}
Hierarchical Gaussian enrichment and overlap identities produce a cascade description. Contact-point supersolution estimates for the pattern-density evolution supply the nonconvex upper bound, while cavity calculations give the matching lower bound and approximation treats merely Borel potentials.

\subsection{Scope relative to this problem}
The Ising Gaussian perceptron is claimed to have an exact pressure formula for every bounded Borel log-potential at positive temperature and density. The question retains these same parameters and conventions. Within these assumptions, this is a claimed affirmative answer to the stated question.

\subsection{Formal verification}
For family 222, the spherical Gaussian perceptron pressure formula is included. The Ising companion entry only proves finite variational value for bounded continuous log-potentials; it neither gives its pressure convergence nor the bounded-Borel extension. A spherical-field supplement supplies a dual approximation formula. The listed Comparator files are \texttt{IsingFiniteness.lean}, \texttt{PerceptronFreeEnergy.lean}, \texttt{SphericalField.lean}. The inventory does not supply a named declaration for these Comparator entries. This site has not executed the Lean code. The scope descriptions do not certify the remaining manuscript arguments.

\subsection{Status caveat}
OpenAI attributes these manuscripts to work by an internal model. Its release cautions that arguments without formalization can contain mistakes. This catalogue has not independently checked the proofs; the labels report the mathematical claims attributed to OpenAI.

\subsection{Sources}
\begin{itemize}
\item OpenAI, \emph{The free energy of the Ising random perceptron}, 2026-09-24; citation key \texttt{OAI:The-free-energy-of-the-Ising-random-perceptron-September-24-2026}. \href{https://github.com/openai/math/blob/adc7f1241b42e322a6451854ab7e4b4c146bf78a/preprints/The-free-energy-of-the-Ising-random-perceptron-September-24-2026/The-free-energy-of-the-Ising-random-perceptron-September-24-2026.pdf}{PDF (GitHub)}; \href{https://github.com/openai/math/raw/adc7f1241b42e322a6451854ab7e4b4c146bf78a/preprints/The-free-energy-of-the-Ising-random-perceptron-September-24-2026/The-free-energy-of-the-Ising-random-perceptron-September-24-2026.pdf}{PDF (direct)}; \href{https://github.com/openai/math/blob/adc7f1241b42e322a6451854ab7e4b4c146bf78a/preprints/The-free-energy-of-the-Ising-random-perceptron-September-24-2026/README.md}{Manuscript page and citation}; \href{https://github.com/openai/math/blob/main/preprints/The-free-energy-of-the-Ising-random-perceptron-September-24-2026/The-free-energy-of-the-Ising-random-perceptron-September-24-2026.pdf}{Latest version}.
\item Family 222: \href{https://github.com/openai/math/blob/adc7f1241b42e322a6451854ab7e4b4c146bf78a/lean/docs/222.md}{Lean scope}.
\item \href{https://github.com/openai/math/blob/adc7f1241b42e322a6451854ab7e4b4c146bf78a/lean/ComparatorChallenges/IsingFiniteness.json}{Comparator IsingFiniteness.json}.
\item \href{https://github.com/openai/math/blob/adc7f1241b42e322a6451854ab7e4b4c146bf78a/lean/ComparatorChallenges/IsingFiniteness.lean}{Comparator IsingFiniteness.lean}.
\item \href{https://github.com/openai/math/blob/adc7f1241b42e322a6451854ab7e4b4c146bf78a/lean/ComparatorChallenges/PerceptronFreeEnergy.json}{Comparator PerceptronFreeEnergy.json}.
\item \href{https://github.com/openai/math/blob/adc7f1241b42e322a6451854ab7e4b4c146bf78a/lean/ComparatorChallenges/PerceptronFreeEnergy.lean}{Comparator PerceptronFreeEnergy.lean}.
\item \href{https://github.com/openai/math/blob/adc7f1241b42e322a6451854ab7e4b4c146bf78a/lean/ComparatorChallenges/SphericalField.json}{Comparator SphericalField.json}.
\item \href{https://github.com/openai/math/blob/adc7f1241b42e322a6451854ab7e4b4c146bf78a/lean/ComparatorChallenges/SphericalField.lean}{Comparator SphericalField.lean}.
\item \href{https://github.com/openai/math/blob/adc7f1241b42e322a6451854ab7e4b4c146bf78a/CONTENTS.md}{OpenAI catalogue entry}.
\end{itemize}
\end{document}
